\documentclass[11pt]{article}
\usepackage{amsmath}
\usepackage[hidelinks]{hyperref}
\begin{document}
\title{A native joint cave-planning comparison at a coupled recovery cut}
\author{}
\date{26 November 2026}
\maketitle

\paragraph{Audit.}
Known travel blockage and a loss of direct radio contact affect the
same outstanding report, not just two unrelated parts of a mission.
The joint combined case can be recovered by a mule, but its position
at the early branch cut matters. Two known-world feasible plans
need not yield a feasible locally selected plan. With one common
staging cell and one timely qualified local cue, the declared
positive case is feasible. The native COPS model and an independent
state search check the required subplans; the equally informed direct
controller and the shared checker accept the same policy. This is a
reduction/calibration result, not an advantage or I2 acceptance.

\paragraph{Baseline and adapters.}
\href{https://ntrs.nasa.gov/citations/20220001447}{Klaesson et al.'s
cave-planning work} separates movement from information flow,
allows carried data and explicit communication, and uses master-plan
flow to restrict the movements of agents that do not know a plan.
Its TRACE executive implements coordination under delays; no
generic global-stop assumption is attributed to it. Here the
actual pinned \href{https://github.com/FilipKlaesson/cops/tree/67d2e5dd4534b3f3eec95b6cfda9d4c9c1746ef0}{COPS}
\texttt{solve\_flow} formulation, indices, objective and plan decoder
run unchanged on a small ordered graph-shaped adapter. The numerical
solver backend is HiGHS rather than the upstream solver options. It
uses the original variable bounds and integrality and checks returned
rows before interpreting a solution. Unused imports for graph
visualization, coloring and clustering are inert, and unsupported
operations fail. This is not a patched-by-hand approximation of the
master constraints, nor the full upstream installation. Extra
time-indexed constraints for the authored inspection, exit sites,
contact changes and optional fixed prefix are explicit. The
independent reference enumerates directed move and contact transitions
with carried report and master knowledge, matching the modeled
instantaneous same-slot communication closure. It is not COPS
under a different label.

\paragraph{Civil fixture.}
One scout starts at S holding the designated report and has to
inspect I at tick 3 and reach E at 4. A mobile helper starts at M,
and a passive base stays at B. Lead motion, relay motion and
permitted communication are distinct typed directed edges.
The scout's normal shortcut reaches N at tick 1 and can use direct
contact N--B. A blockage removes S--N. The scout can instead travel
S--X--Y--I--E; without partition, X--B still gives the report to
base at 1. If only the direct links fail, the reachable N--M
rendezvous passes custody to the helper and its movement brings the
report to a checkpoint Z with Z--B contact at 3. When both changes
occur, the shortcut rendezvous is no longer reachable and X--B
is lost. The same old report has to use Y--U or Y--G as the lead
passes Y at 2. The helper later reaches Z and base gets it at 3.
The native known-graph model admits a two-hop route for each possible
remaining contact. The helper is required to end safely (M or Z);
weather, traversal robustness, transfer durations finer than a tick,
and report delivery beyond the tested terminal receiver are not
certified. A packet delivery is not a robot position claim.

In the compound case, an additional contact change removes one
of Y--U and Y--G before handoff. Which remains valid is not
available before the helper's first step. There are no hidden
poses or branch identifiers at that cut. The fault envelope
represents known initial blockage and partition and a local sensor
law at time 1. The solver used with the full graph per case is
an oracle about the hidden branch until that law supplies the cue.
Accounting separately for that evidence is essential: successful
plans synthesized for the two full graphs cannot both be supplied
at time zero as locally selected outcomes.

\paragraph{Two cuts and executable guard.}
Without the optional common staging cell J, the first-step options
for U and G are disjoint: for U the helper can stay at M or move
to U immediately, while for G it must start toward W at tick 0.
A received cue after that step does not repair the wrong committed
route. Count this as a physical prefix obstruction, not a deficit
in the COPS solver. The common staging move M--J makes both possible,
with J--U and J--G each one step. At tick 1 it must learn which
to choose for the appointment. The native decoded witnesses
contain both actual communication edges at 2 and 3, the scout's
continued inspection and exit, and a helper egress. The common prefix is not free: with branch
already known the native optimum uses two helper moves for U and
three for G, compared to three for either staged branch. If the cue
is missing, a preauthorized retreat J--M still exists, and
inspection and exits are possible without meeting the old report
deadline; that is a constrained fallback, not successful service.
If the helper did not install the plan and no authorized channel
reaches it before the first move, the original master constraints
reject the first staging move. Preinstallation is supplied in
the positive case to both controllers, not smuggled in as a free
late master signal.

A separate staged pinned native SubT broker control sends a one-byte
branch cue to the helper endpoint, checks its own receive callback,
and takes the second-route action based on that receiver view.
The lead-to-helper data handoff and helper-to-base packet are only
counted after actual receipt. Transmit counters charge two bytes
per attempted report hop, including failed attempts, and both
controllers use the same one-byte cue attempt. The input geometry chooses which radio
leg is present; it is not inferred from a packet. In these tests the
broker actually delivers the short cue in each qualified positive
path; the underlying truth, ownership and validity horizon of the
cue are authored. Erased, late and wrong-version receipts do not
supply a branch. A delivered cue followed by another change can
be wrong about the still-available route; the initial receipt does
not certify a support law through the handoff. The forced invalidated
cue is a negative control, not a claimed supported policy.
An unrelated native callback at 4 is preserved even on the
failed report paths.

\paragraph{Results and comparison boundary.}
The four factorial classes separate nominal, blockage-only,
partition-only and joint changes. Earliest admitted base receipt
in the native plan is 1, 1, 3 and 3, respectively, but the
last requires a new reachable transfer and not the old rendezvous.
There are 55 matched COPS/direct conditions, including master
scope and forced prefix/branch tests, 120 native packet controls
and ten same-process local selections. Two abstract finite cut
lemmas were checked in Lean. Exact schedules are fixture-relative;
observational slots for cue receipt are charged and not equated with
learning a real mine map. The result shows an established flow
formulation can solve the known subproblems and a standard common
local-prefix search can assemble the tested adaptive policy; no
claim is made that the source package already solves the larger
partial-information fork in a single call.
No end-to-end comparison to TRACE, ACHORD's full live
\href{https://lillyclark.github.io/files/ACHORD_RA_L_IROS_.pdf}{networking
and return behavior}, or reproducible field rate is claimed.
A valid continuation experiment should now test the case where
choosing a branch changes an already staged old effect or the
future opportunity to obtain evidence, without giving the adapted
joint planner less authority or information.
\end{document}
